\section{The causal completion problem}\label{sec:problem}
A prepared experiment determines probabilities of executable future tests. It does not endow its observer with a real-valued posterior register, a free clock, or a chosen manifold. The foundational question is how actual acquisition and permissible causal processing determine the resolution of a finite observer and the risk of its decisions. We retain the order
\[
\begin{gathered}
\text{prepared kernels}\longrightarrow\text{future tests}\longrightarrow
\text{predictive quotient}\\
\longrightarrow\text{acquired law}\longrightarrow
\text{causal implementation and risk}.
\end{gathered}
\]
Statistical geometry is a derived object in this chain. A formally reachable set is not an acquisition distribution, and terminal quantization need not describe an executable online observer.

The central obstruction considered here has two parts. Information compressed at an intermediate cut must still support the next experiment, and the finite machine must know its position in the protocol. Conditional-barycenter pooling addresses the first part: each retained label is calibrated under the law of the implemented observer, rather than assigned an approximation to an inaccessible filter. Its loss is an exact Jensen defect of the optimal continuation value. The second part cannot in general be removed by reusing labels at different phases. When all report laws have a common positive support, an exact internal deadline forces distinct phase supports.

Theorem~\ref{thm:main} joins these mechanisms. Its generality is a finite affine predictive class derived from positive experimental instruments; it is not limited to a probability simplex or a commuting measurement algebra. The hypotheses are uniform statements about each raw one-step acquired law and its stability, not properties of a selected optimal code. The decisive improvement is the joint law
\begin{equation}\label{eq:intro}
 R_\cp(n,M)-B_n^*\asymp M^{-2/d},\qquad
 R_\aut(n,M)-B_n^*\asymp(M-n)^{-2/d},\quad M\ge n+1.
\end{equation}
Here $B_n^*$ is the optimal full-history risk in the same $n$-acquisition experiment. The constants are independent of $n$ and $M$. An autonomous class with $M\le n$ is empty under the exact-deadline convention. Thus horizon-uniform completion does not mean ignoring the clock: for example, uniform comparison at the same budget follows when $M\ge(1+\eta)n$, with a constant depending on $\eta>0$.

The proof assigns few labels to phases whose continuation values have little sensitivity and more labels to the final phases. Block contraction suffices; individual one-step coefficients may exceed one. A convex extension argument permits acquired laws to approach the boundary, which is essential for ordinary unbounded Gaussian reports. Two structurally different realizations follow: classical hidden-state sensing and noncommuting positive quantum instruments. Their predictive dimensions are obtained from executable tests and raw changes of variables. Neither realization is used to prove the general theorem. The singular extension treats actual mixtures of atoms and continuous laws, preserving their subprobability masses.

The posterior-pooling identity and its interior curvature argument are inherited from the preceding article, supplied intact as Supplement U~\cite{pooling}. The complete separated-refinement and continuation-cut results remain in Supplement T~\cite{retained}, and the block-filtering and stopped-task proofs remain in Supplement S~\cite{companion}. The new claims here are the boundary-valid transfer, the sharp charged-clock law, its block-stable task-general proof, and the Gaussian and noncommuting-instrument verifications. The state-space language itself is classical~\cite{davies,blackwell}; it is not claimed as new mathematics.

\subsection{Prepared kernels and affine predictive realizations}
All measurable spaces are standard Borel. Preparation and legal interventions specify positive normalized kernels $K_t^\lambda(dx',dy,dc\mid x,a)$. Failed observations, physical duration and intervention costs belong to $c$. Policies use only past reports, actions and the declared interface. They cannot read the unknown parameter $\lambda$. The rate theorems below fix a known prepared experiment; a Bayes statement for that preparation is not a learning theorem or a simultaneous sufficiency theorem for every parameter.

An executable test is a legal finite continuation with a bounded scored readout. Two histories at a phase are equivalent when they have the same legal interface and the same expectation for every such test. Appendix~\ref{sec:quotient} gives measurable hypotheses for this quotient. The finite-dimensional hypothesis used in the main theorem is the following additional, verifiable closure property.

\begin{definition}[Affine predictive realization]\label{def:affine}
The experiment has a compact convex predictive carrier $K$ of affine dimension $d\ge1$, with a fixed affine coordinate identification $K\subset\R^d$. Every relevant future-test expectation is affine on $K$, and the tests separate its points. For every phase $t$ and legal action $a$, the unnormalized report-and-continuation functionals depend affinely on the incoming state. Their normalization gives a jointly Borel report density $r_t(s,a,y)$ and updated state $F_t(s,a,y)\in K$, relative to a fixed report measure $m^a$. They realize the actual report and future-test probabilities. Mixtures of preparations have the corresponding barycentric state.

The conditions hold on all $K$, including mixtures created by compression, not only on the image of a preferred full-history policy. Write
\begin{equation}\label{eq:Lambda}
 J_t(s,a)(dy)=r_t(s,a,y)m^a(dy),\qquad
 \Lambda_t(s,a)=\law_{Y\sim J_t(s,a)}F_t(s,a,Y).
\end{equation}
\end{definition}
One way to verify the definition is to exhibit a finite-dimensional span of executable tests, containing the constant test and invariant under the unnormalized report instruments, then evaluate a basis on histories and mixtures. Positivity and normalization descend from the original experiment. Neither covariance nor Fisher information is an input. Closure and the existence of the jointly defined update are assumptions to verify, not consequences of finite rank alone. The actual acquired law at a cut is always the pushforward of the specified preparation and policy, never an arbitrary measure on $K$.

We fix a norm $\norm{\cdot}$ on the translation space of $K$. Wasserstein distance $W_1$ uses this norm. Lebesgue density and strong concavity below use the displayed Euclidean coordinates; equivalence constants of these norms are fixed experimental constants. Affine coordinate changes alter the constants, not the asserted exponents or state counts.

\subsection{Task and autonomous state model}
There is a common nonempty finite action alphabet $\mathcal A$ at every phase. After exactly $n$ physical calls, an irreversible terminal decision is scored. Its conditional expected loss $\ell(s,u)\in[0,H]$ is affine in $s$ for each decision $u$; a Borel minimizing decision exists. Put
\begin{equation}\label{eq:g}
 g(s)=\min_u\ell(s,u),\qquad V_n=g,\qquad
 V_t(s)=\min_{a\in\mathcal A}\int V_{t+1}(z)\Lambda_t(s,a)(dz),\quad B_n^*=V_0(s_0).
\end{equation}
The affine risks imply concavity. We assume $g$ is $L_g$-Lipschitz and $\alpha$-strongly concave in Euclidean coordinates, with $\alpha>0$. The latter is needed only for the quantitative converse. Upper transfer and the exact Jensen identity remain valid without it. These assumptions cover more than quadratic losses, as shown after the main theorem.

An autonomous observer has at most $M$ persistent labels. Its action, whether to stop, and terminal decision depend on the current label alone. The next label may use the current report and fresh independent randomization. At persistent cuts there is no readable elapsed time, old output tape, history, instruction pointer or persistent random seed outside this label. Temporary computation is erased. A fixed read-only programme may depend on the known kernels, preparation, horizon and budgets, but not on acquired data or an unknown parameter. Exact real instructions belong to the cardinality model only; Section~\ref{sec:resources} distinguishes finite-bit execution. The observer must itself stop after exactly $n$ calls almost surely. A terminal request supplied by an external clock is a different resource model.

A checkpoint observer may use full history during acquisition and keeps only $M$ terminal labels. Both classes optimize over all admitted full-history sensing policies, use the same preparation and scored task, and pay the same physical call cost. For a deterministic policy $\pi$, write $\nu_n^\pi=\law(S_n^\pi)$ and $B_n^\pi=\E g(S_n^\pi)$. Define the task distortion
\begin{equation}\label{eq:taskdist}
 D_M^g(\nu)=\inf_{\mathcal P:\,|\mathcal P|\le M}
 \sum_{C\in\mathcal P}\left[\nu(C)g\!\left(\frac{\int_Cs\,\nu(ds)}{\nu(C)}\right)-\int_C g(s)\nu(ds)\right].
\end{equation}
Zero-mass cells are omitted. This is a distortion of the actual law. Independent seed conditioning gives
\begin{equation}\label{eq:cpvar}
 R_\cp(n,M)=\inf_\pi\{B_n^\pi+D_M^g(\nu_n^\pi)\}.
\end{equation}
For fixed terminal readouts, assigning a history to a least-loss readout is a Borel function of its predictive state: each conditional readout risk is affine in that state. Replacing any history-based randomized encoder by this assignment cannot increase risk, and then optimizing each cell readout gives the barycentric distortion above. Conditioning independent policy seeds completes the reduction; randomized decisions do not improve a linear conditional loss. This is a checkpoint relaxation, not purification of a clock-free autonomous machine. Policy-dependent summands in \eqref{eq:cpvar} are not minimized separately.
